% ============================================================================
% ESQUEMA PROPOSTO (banca, sugestão — não mergeado): o gate do oráculo e a
% régua da hipótese como fluxo de decisão. Complementa a Seção
% sec:metodo-oraculo-decisao e a Seção sec:intro-hipotese.
%
% Uso no Cap. 3 (dentro de um ambiente figure):
%   % ============================================================================
% ESQUEMA PROPOSTO (banca, sugestão — não mergeado): o gate do oráculo e a
% régua da hipótese como fluxo de decisão. Complementa a Seção
% sec:metodo-oraculo-decisao e a Seção sec:intro-hipotese.
%
% Uso no Cap. 3 (dentro de um ambiente figure):
%   % ============================================================================
% ESQUEMA PROPOSTO (banca, sugestão — não mergeado): o gate do oráculo e a
% régua da hipótese como fluxo de decisão. Complementa a Seção
% sec:metodo-oraculo-decisao e a Seção sec:intro-hipotese.
%
% Uso no Cap. 3 (dentro de um ambiente figure):
%   % ============================================================================
% ESQUEMA PROPOSTO (banca, sugestão — não mergeado): o gate do oráculo e a
% régua da hipótese como fluxo de decisão. Complementa a Seção
% sec:metodo-oraculo-decisao e a Seção sec:intro-hipotese.
%
% Uso no Cap. 3 (dentro de um ambiente figure):
%   \input{3-metodo/esquemas-propostos/esq-gate-e-regua}
% Uso standalone (pré-visualização):
%   \documentclass[tikz,border=6pt]{standalone}
%   \usetikzlibrary{arrows.meta, positioning}
%   \begin{document}\input{esq-gate-e-regua.tex}\end{document}
%   (no modo standalone, troque as \ref{...} por texto fixo ou defina-as)
% Requer apenas as bibliotecas já carregadas em packages.tex
% (arrows.meta, positioning). Desenhado para ser legível em preto e branco:
% decisão = borda grossa + cinza escuro; medição = cinza claro;
% consequência = cinza médio; condicional/remissão = tracejado.
% FREEZE respeitado: todos os números são os já fixados no texto
% (85%; 0,95 x 89,56%; 0,95 x acc(D); 34.724; 231.490; 15%).
% Compilado e inspecionado visualmente (pdflatex + render PNG) — ver
% NOTA-esquemas.md, iterações 4-5 do loop.
% ============================================================================
\begin{tikzpicture}[
    x=1cm, y=1cm,
    caixa/.style={draw, rounded corners=2pt, align=center, text width=34mm,
                  minimum height=11mm, inner sep=3pt, font=\footnotesize},
    decisao/.style={caixa, thick, fill=gray!25},
    medida/.style={caixa, fill=gray!8},
    consequencia/.style={caixa, fill=gray!14},
    seta/.style={-{Stealth[length=2.4mm]}, thick},
    rot/.style={font=\scriptsize, align=center, inner sep=2pt}
]
    % ---------------- faixa superior: o GATE do oráculo -------------------
    \node[rot, anchor=west] (tit1) at (-1.9,1.7)
        {\textbf{Gate do oráculo} (fixado de antemão, Seção~\ref{sec:metodo-oraculo-decisao})};

    \node[medida]  (e0)   at (0,0)
        {avaliação dos oráculos\\LLM (S-rand $n=1.000$;\\S-strat $n\approx1.863$)};
    \node[decisao] (gate) at (4.6,0)
        {algum oráculo com\\acurácia $\ge 85\%$\\na S-rand?};
    \node[caixa]   (sim)  at (9.4,1.3)
        {oráculo \textbf{aprovado}:\\papéis pelo critério\\original};
    \node[consequencia] (nao) at (9.4,-1.3)
        {robustez ao ruído\\\textbf{obrigatória}; oráculo ruidoso\\vira cenário central};

    \draw[seta] (e0) -- (gate);
    % âncoras deixam o rótulo inteiro do lado livre da diagonal
    % (a seta cortava o "m"/"o" — achado da inspeção visual)
    \draw[seta] (gate.east) -- node[rot, anchor=south east, pos=0.55] {sim} (sim.west);
    \draw[seta] (gate.east) -- node[rot, anchor=north east, pos=0.55] {não} (nao.west);

    \node[consequencia, text width=44mm] (papeis) at (9.4,-3.5)
        {papéis: LLM Inicial pela razão acurácia/custo; LLM Avançado se
         superior com significância (senão a Fase~3 é eliminada)};
    \draw[seta] (nao) -- (papeis);

    \node[rot, text width=50mm, anchor=north] (ramo) at (9.4,-4.75)
        {o ramo efetivamente seguido, e a divergência que ele exige
         declarar, são reportados com os resultados (Seção~\ref{sec:res-gate})};
    \draw[seta, dashed] (ramo.north) -- (papeis.south);

    % ---------------- faixa inferior: a RÉGUA da hipótese ------------------
    \node[rot, anchor=west] (tit2) at (-1.9,-5.9)
        {\textbf{Critério da hipótese} (Seção~\ref{sec:intro-hipotese})};

    \node[medida] (regua) at (0,-7.4)
        {régua $D$: o mesmo\\classificador com o\\\textit{pool} de referência\\inteiramente rotulado};
    \node[decisao] (crit) at (4.6,-7.4)
        {$\mathrm{acc} \ge 0{,}95 \times \mathrm{acc}(D)$\\com no máximo\\34.724 rótulos?};
    \node[caixa] (veredito) at (9.4,-7.4)
        {veredito da hipótese\\(julgado nos resultados,\\Seção~\ref{sec:res-e3p})};

    \draw[seta] (regua) -- (crit);
    \draw[seta] (crit) -- (veredito);
    \node[rot, below=2.5mm of crit, text width=62mm]
        {teto de 34.724 rótulos $= 15\%$ da população deduplicada de 231.490;
         a acurácia é a métrica do critério (pré-registro); o Macro F1
         acompanha como análise de robustez};

    % ---------------- amarração entre as duas réguas -----------------------
    % conector vertical reto entre as duas caixas de decisão, com o rótulo
    % ancorado à esquerda da linha (nada cruza título nem seta).
    % Fica por último porque referencia nós das duas faixas.
    \draw[seta, dashed] (gate.south) -- (crit.north);
    \node[rot, text width=40mm, anchor=east] at (4.3,-3.3)
        {a mesma razão de $0{,}95$ nas duas réguas:
         $85\% \approx 0{,}95 \times 89{,}56\%$ do melhor
         baseline supervisionado leve};
\end{tikzpicture}

% Uso standalone (pré-visualização):
%   \documentclass[tikz,border=6pt]{standalone}
%   \usetikzlibrary{arrows.meta, positioning}
%   \begin{document}% ============================================================================
% ESQUEMA PROPOSTO (banca, sugestão — não mergeado): o gate do oráculo e a
% régua da hipótese como fluxo de decisão. Complementa a Seção
% sec:metodo-oraculo-decisao e a Seção sec:intro-hipotese.
%
% Uso no Cap. 3 (dentro de um ambiente figure):
%   \input{3-metodo/esquemas-propostos/esq-gate-e-regua}
% Uso standalone (pré-visualização):
%   \documentclass[tikz,border=6pt]{standalone}
%   \usetikzlibrary{arrows.meta, positioning}
%   \begin{document}\input{esq-gate-e-regua.tex}\end{document}
%   (no modo standalone, troque as \ref{...} por texto fixo ou defina-as)
% Requer apenas as bibliotecas já carregadas em packages.tex
% (arrows.meta, positioning). Desenhado para ser legível em preto e branco:
% decisão = borda grossa + cinza escuro; medição = cinza claro;
% consequência = cinza médio; condicional/remissão = tracejado.
% FREEZE respeitado: todos os números são os já fixados no texto
% (85%; 0,95 x 89,56%; 0,95 x acc(D); 34.724; 231.490; 15%).
% Compilado e inspecionado visualmente (pdflatex + render PNG) — ver
% NOTA-esquemas.md, iterações 4-5 do loop.
% ============================================================================
\begin{tikzpicture}[
    x=1cm, y=1cm,
    caixa/.style={draw, rounded corners=2pt, align=center, text width=34mm,
                  minimum height=11mm, inner sep=3pt, font=\footnotesize},
    decisao/.style={caixa, thick, fill=gray!25},
    medida/.style={caixa, fill=gray!8},
    consequencia/.style={caixa, fill=gray!14},
    seta/.style={-{Stealth[length=2.4mm]}, thick},
    rot/.style={font=\scriptsize, align=center, inner sep=2pt}
]
    % ---------------- faixa superior: o GATE do oráculo -------------------
    \node[rot, anchor=west] (tit1) at (-1.9,1.7)
        {\textbf{Gate do oráculo} (fixado de antemão, Seção~\ref{sec:metodo-oraculo-decisao})};

    \node[medida]  (e0)   at (0,0)
        {avaliação dos oráculos\\LLM (S-rand $n=1.000$;\\S-strat $n\approx1.863$)};
    \node[decisao] (gate) at (4.6,0)
        {algum oráculo com\\acurácia $\ge 85\%$\\na S-rand?};
    \node[caixa]   (sim)  at (9.4,1.3)
        {oráculo \textbf{aprovado}:\\papéis pelo critério\\original};
    \node[consequencia] (nao) at (9.4,-1.3)
        {robustez ao ruído\\\textbf{obrigatória}; oráculo ruidoso\\vira cenário central};

    \draw[seta] (e0) -- (gate);
    % âncoras deixam o rótulo inteiro do lado livre da diagonal
    % (a seta cortava o "m"/"o" — achado da inspeção visual)
    \draw[seta] (gate.east) -- node[rot, anchor=south east, pos=0.55] {sim} (sim.west);
    \draw[seta] (gate.east) -- node[rot, anchor=north east, pos=0.55] {não} (nao.west);

    \node[consequencia, text width=44mm] (papeis) at (9.4,-3.5)
        {papéis: LLM Inicial pela razão acurácia/custo; LLM Avançado se
         superior com significância (senão a Fase~3 é eliminada)};
    \draw[seta] (nao) -- (papeis);

    \node[rot, text width=50mm, anchor=north] (ramo) at (9.4,-4.75)
        {o ramo efetivamente seguido, e a divergência que ele exige
         declarar, são reportados com os resultados (Seção~\ref{sec:res-gate})};
    \draw[seta, dashed] (ramo.north) -- (papeis.south);

    % ---------------- faixa inferior: a RÉGUA da hipótese ------------------
    \node[rot, anchor=west] (tit2) at (-1.9,-5.9)
        {\textbf{Critério da hipótese} (Seção~\ref{sec:intro-hipotese})};

    \node[medida] (regua) at (0,-7.4)
        {régua $D$: o mesmo\\classificador com o\\\textit{pool} de referência\\inteiramente rotulado};
    \node[decisao] (crit) at (4.6,-7.4)
        {$\mathrm{acc} \ge 0{,}95 \times \mathrm{acc}(D)$\\com no máximo\\34.724 rótulos?};
    \node[caixa] (veredito) at (9.4,-7.4)
        {veredito da hipótese\\(julgado nos resultados,\\Seção~\ref{sec:res-e3p})};

    \draw[seta] (regua) -- (crit);
    \draw[seta] (crit) -- (veredito);
    \node[rot, below=2.5mm of crit, text width=62mm]
        {teto de 34.724 rótulos $= 15\%$ da população deduplicada de 231.490;
         a acurácia é a métrica do critério (pré-registro); o Macro F1
         acompanha como análise de robustez};

    % ---------------- amarração entre as duas réguas -----------------------
    % conector vertical reto entre as duas caixas de decisão, com o rótulo
    % ancorado à esquerda da linha (nada cruza título nem seta).
    % Fica por último porque referencia nós das duas faixas.
    \draw[seta, dashed] (gate.south) -- (crit.north);
    \node[rot, text width=40mm, anchor=east] at (4.3,-3.3)
        {a mesma razão de $0{,}95$ nas duas réguas:
         $85\% \approx 0{,}95 \times 89{,}56\%$ do melhor
         baseline supervisionado leve};
\end{tikzpicture}
\end{document}
%   (no modo standalone, troque as \ref{...} por texto fixo ou defina-as)
% Requer apenas as bibliotecas já carregadas em packages.tex
% (arrows.meta, positioning). Desenhado para ser legível em preto e branco:
% decisão = borda grossa + cinza escuro; medição = cinza claro;
% consequência = cinza médio; condicional/remissão = tracejado.
% FREEZE respeitado: todos os números são os já fixados no texto
% (85%; 0,95 x 89,56%; 0,95 x acc(D); 34.724; 231.490; 15%).
% Compilado e inspecionado visualmente (pdflatex + render PNG) — ver
% NOTA-esquemas.md, iterações 4-5 do loop.
% ============================================================================
\begin{tikzpicture}[
    x=1cm, y=1cm,
    caixa/.style={draw, rounded corners=2pt, align=center, text width=34mm,
                  minimum height=11mm, inner sep=3pt, font=\footnotesize},
    decisao/.style={caixa, thick, fill=gray!25},
    medida/.style={caixa, fill=gray!8},
    consequencia/.style={caixa, fill=gray!14},
    seta/.style={-{Stealth[length=2.4mm]}, thick},
    rot/.style={font=\scriptsize, align=center, inner sep=2pt}
]
    % ---------------- faixa superior: o GATE do oráculo -------------------
    \node[rot, anchor=west] (tit1) at (-1.9,1.7)
        {\textbf{Gate do oráculo} (fixado de antemão, Seção~\ref{sec:metodo-oraculo-decisao})};

    \node[medida]  (e0)   at (0,0)
        {avaliação dos oráculos\\LLM (S-rand $n=1.000$;\\S-strat $n\approx1.863$)};
    \node[decisao] (gate) at (4.6,0)
        {algum oráculo com\\acurácia $\ge 85\%$\\na S-rand?};
    \node[caixa]   (sim)  at (9.4,1.3)
        {oráculo \textbf{aprovado}:\\papéis pelo critério\\original};
    \node[consequencia] (nao) at (9.4,-1.3)
        {robustez ao ruído\\\textbf{obrigatória}; oráculo ruidoso\\vira cenário central};

    \draw[seta] (e0) -- (gate);
    % âncoras deixam o rótulo inteiro do lado livre da diagonal
    % (a seta cortava o "m"/"o" — achado da inspeção visual)
    \draw[seta] (gate.east) -- node[rot, anchor=south east, pos=0.55] {sim} (sim.west);
    \draw[seta] (gate.east) -- node[rot, anchor=north east, pos=0.55] {não} (nao.west);

    \node[consequencia, text width=44mm] (papeis) at (9.4,-3.5)
        {papéis: LLM Inicial pela razão acurácia/custo; LLM Avançado se
         superior com significância (senão a Fase~3 é eliminada)};
    \draw[seta] (nao) -- (papeis);

    \node[rot, text width=50mm, anchor=north] (ramo) at (9.4,-4.75)
        {o ramo efetivamente seguido, e a divergência que ele exige
         declarar, são reportados com os resultados (Seção~\ref{sec:res-gate})};
    \draw[seta, dashed] (ramo.north) -- (papeis.south);

    % ---------------- faixa inferior: a RÉGUA da hipótese ------------------
    \node[rot, anchor=west] (tit2) at (-1.9,-5.9)
        {\textbf{Critério da hipótese} (Seção~\ref{sec:intro-hipotese})};

    \node[medida] (regua) at (0,-7.4)
        {régua $D$: o mesmo\\classificador com o\\\textit{pool} de referência\\inteiramente rotulado};
    \node[decisao] (crit) at (4.6,-7.4)
        {$\mathrm{acc} \ge 0{,}95 \times \mathrm{acc}(D)$\\com no máximo\\34.724 rótulos?};
    \node[caixa] (veredito) at (9.4,-7.4)
        {veredito da hipótese\\(julgado nos resultados,\\Seção~\ref{sec:res-e3p})};

    \draw[seta] (regua) -- (crit);
    \draw[seta] (crit) -- (veredito);
    \node[rot, below=2.5mm of crit, text width=62mm]
        {teto de 34.724 rótulos $= 15\%$ da população deduplicada de 231.490;
         a acurácia é a métrica do critério (pré-registro); o Macro F1
         acompanha como análise de robustez};

    % ---------------- amarração entre as duas réguas -----------------------
    % conector vertical reto entre as duas caixas de decisão, com o rótulo
    % ancorado à esquerda da linha (nada cruza título nem seta).
    % Fica por último porque referencia nós das duas faixas.
    \draw[seta, dashed] (gate.south) -- (crit.north);
    \node[rot, text width=40mm, anchor=east] at (4.3,-3.3)
        {a mesma razão de $0{,}95$ nas duas réguas:
         $85\% \approx 0{,}95 \times 89{,}56\%$ do melhor
         baseline supervisionado leve};
\end{tikzpicture}

% Uso standalone (pré-visualização):
%   \documentclass[tikz,border=6pt]{standalone}
%   \usetikzlibrary{arrows.meta, positioning}
%   \begin{document}% ============================================================================
% ESQUEMA PROPOSTO (banca, sugestão — não mergeado): o gate do oráculo e a
% régua da hipótese como fluxo de decisão. Complementa a Seção
% sec:metodo-oraculo-decisao e a Seção sec:intro-hipotese.
%
% Uso no Cap. 3 (dentro de um ambiente figure):
%   % ============================================================================
% ESQUEMA PROPOSTO (banca, sugestão — não mergeado): o gate do oráculo e a
% régua da hipótese como fluxo de decisão. Complementa a Seção
% sec:metodo-oraculo-decisao e a Seção sec:intro-hipotese.
%
% Uso no Cap. 3 (dentro de um ambiente figure):
%   \input{3-metodo/esquemas-propostos/esq-gate-e-regua}
% Uso standalone (pré-visualização):
%   \documentclass[tikz,border=6pt]{standalone}
%   \usetikzlibrary{arrows.meta, positioning}
%   \begin{document}\input{esq-gate-e-regua.tex}\end{document}
%   (no modo standalone, troque as \ref{...} por texto fixo ou defina-as)
% Requer apenas as bibliotecas já carregadas em packages.tex
% (arrows.meta, positioning). Desenhado para ser legível em preto e branco:
% decisão = borda grossa + cinza escuro; medição = cinza claro;
% consequência = cinza médio; condicional/remissão = tracejado.
% FREEZE respeitado: todos os números são os já fixados no texto
% (85%; 0,95 x 89,56%; 0,95 x acc(D); 34.724; 231.490; 15%).
% Compilado e inspecionado visualmente (pdflatex + render PNG) — ver
% NOTA-esquemas.md, iterações 4-5 do loop.
% ============================================================================
\begin{tikzpicture}[
    x=1cm, y=1cm,
    caixa/.style={draw, rounded corners=2pt, align=center, text width=34mm,
                  minimum height=11mm, inner sep=3pt, font=\footnotesize},
    decisao/.style={caixa, thick, fill=gray!25},
    medida/.style={caixa, fill=gray!8},
    consequencia/.style={caixa, fill=gray!14},
    seta/.style={-{Stealth[length=2.4mm]}, thick},
    rot/.style={font=\scriptsize, align=center, inner sep=2pt}
]
    % ---------------- faixa superior: o GATE do oráculo -------------------
    \node[rot, anchor=west] (tit1) at (-1.9,1.7)
        {\textbf{Gate do oráculo} (fixado de antemão, Seção~\ref{sec:metodo-oraculo-decisao})};

    \node[medida]  (e0)   at (0,0)
        {avaliação dos oráculos\\LLM (S-rand $n=1.000$;\\S-strat $n\approx1.863$)};
    \node[decisao] (gate) at (4.6,0)
        {algum oráculo com\\acurácia $\ge 85\%$\\na S-rand?};
    \node[caixa]   (sim)  at (9.4,1.3)
        {oráculo \textbf{aprovado}:\\papéis pelo critério\\original};
    \node[consequencia] (nao) at (9.4,-1.3)
        {robustez ao ruído\\\textbf{obrigatória}; oráculo ruidoso\\vira cenário central};

    \draw[seta] (e0) -- (gate);
    % âncoras deixam o rótulo inteiro do lado livre da diagonal
    % (a seta cortava o "m"/"o" — achado da inspeção visual)
    \draw[seta] (gate.east) -- node[rot, anchor=south east, pos=0.55] {sim} (sim.west);
    \draw[seta] (gate.east) -- node[rot, anchor=north east, pos=0.55] {não} (nao.west);

    \node[consequencia, text width=44mm] (papeis) at (9.4,-3.5)
        {papéis: LLM Inicial pela razão acurácia/custo; LLM Avançado se
         superior com significância (senão a Fase~3 é eliminada)};
    \draw[seta] (nao) -- (papeis);

    \node[rot, text width=50mm, anchor=north] (ramo) at (9.4,-4.75)
        {o ramo efetivamente seguido, e a divergência que ele exige
         declarar, são reportados com os resultados (Seção~\ref{sec:res-gate})};
    \draw[seta, dashed] (ramo.north) -- (papeis.south);

    % ---------------- faixa inferior: a RÉGUA da hipótese ------------------
    \node[rot, anchor=west] (tit2) at (-1.9,-5.9)
        {\textbf{Critério da hipótese} (Seção~\ref{sec:intro-hipotese})};

    \node[medida] (regua) at (0,-7.4)
        {régua $D$: o mesmo\\classificador com o\\\textit{pool} de referência\\inteiramente rotulado};
    \node[decisao] (crit) at (4.6,-7.4)
        {$\mathrm{acc} \ge 0{,}95 \times \mathrm{acc}(D)$\\com no máximo\\34.724 rótulos?};
    \node[caixa] (veredito) at (9.4,-7.4)
        {veredito da hipótese\\(julgado nos resultados,\\Seção~\ref{sec:res-e3p})};

    \draw[seta] (regua) -- (crit);
    \draw[seta] (crit) -- (veredito);
    \node[rot, below=2.5mm of crit, text width=62mm]
        {teto de 34.724 rótulos $= 15\%$ da população deduplicada de 231.490;
         a acurácia é a métrica do critério (pré-registro); o Macro F1
         acompanha como análise de robustez};

    % ---------------- amarração entre as duas réguas -----------------------
    % conector vertical reto entre as duas caixas de decisão, com o rótulo
    % ancorado à esquerda da linha (nada cruza título nem seta).
    % Fica por último porque referencia nós das duas faixas.
    \draw[seta, dashed] (gate.south) -- (crit.north);
    \node[rot, text width=40mm, anchor=east] at (4.3,-3.3)
        {a mesma razão de $0{,}95$ nas duas réguas:
         $85\% \approx 0{,}95 \times 89{,}56\%$ do melhor
         baseline supervisionado leve};
\end{tikzpicture}

% Uso standalone (pré-visualização):
%   \documentclass[tikz,border=6pt]{standalone}
%   \usetikzlibrary{arrows.meta, positioning}
%   \begin{document}% ============================================================================
% ESQUEMA PROPOSTO (banca, sugestão — não mergeado): o gate do oráculo e a
% régua da hipótese como fluxo de decisão. Complementa a Seção
% sec:metodo-oraculo-decisao e a Seção sec:intro-hipotese.
%
% Uso no Cap. 3 (dentro de um ambiente figure):
%   \input{3-metodo/esquemas-propostos/esq-gate-e-regua}
% Uso standalone (pré-visualização):
%   \documentclass[tikz,border=6pt]{standalone}
%   \usetikzlibrary{arrows.meta, positioning}
%   \begin{document}\input{esq-gate-e-regua.tex}\end{document}
%   (no modo standalone, troque as \ref{...} por texto fixo ou defina-as)
% Requer apenas as bibliotecas já carregadas em packages.tex
% (arrows.meta, positioning). Desenhado para ser legível em preto e branco:
% decisão = borda grossa + cinza escuro; medição = cinza claro;
% consequência = cinza médio; condicional/remissão = tracejado.
% FREEZE respeitado: todos os números são os já fixados no texto
% (85%; 0,95 x 89,56%; 0,95 x acc(D); 34.724; 231.490; 15%).
% Compilado e inspecionado visualmente (pdflatex + render PNG) — ver
% NOTA-esquemas.md, iterações 4-5 do loop.
% ============================================================================
\begin{tikzpicture}[
    x=1cm, y=1cm,
    caixa/.style={draw, rounded corners=2pt, align=center, text width=34mm,
                  minimum height=11mm, inner sep=3pt, font=\footnotesize},
    decisao/.style={caixa, thick, fill=gray!25},
    medida/.style={caixa, fill=gray!8},
    consequencia/.style={caixa, fill=gray!14},
    seta/.style={-{Stealth[length=2.4mm]}, thick},
    rot/.style={font=\scriptsize, align=center, inner sep=2pt}
]
    % ---------------- faixa superior: o GATE do oráculo -------------------
    \node[rot, anchor=west] (tit1) at (-1.9,1.7)
        {\textbf{Gate do oráculo} (fixado de antemão, Seção~\ref{sec:metodo-oraculo-decisao})};

    \node[medida]  (e0)   at (0,0)
        {avaliação dos oráculos\\LLM (S-rand $n=1.000$;\\S-strat $n\approx1.863$)};
    \node[decisao] (gate) at (4.6,0)
        {algum oráculo com\\acurácia $\ge 85\%$\\na S-rand?};
    \node[caixa]   (sim)  at (9.4,1.3)
        {oráculo \textbf{aprovado}:\\papéis pelo critério\\original};
    \node[consequencia] (nao) at (9.4,-1.3)
        {robustez ao ruído\\\textbf{obrigatória}; oráculo ruidoso\\vira cenário central};

    \draw[seta] (e0) -- (gate);
    % âncoras deixam o rótulo inteiro do lado livre da diagonal
    % (a seta cortava o "m"/"o" — achado da inspeção visual)
    \draw[seta] (gate.east) -- node[rot, anchor=south east, pos=0.55] {sim} (sim.west);
    \draw[seta] (gate.east) -- node[rot, anchor=north east, pos=0.55] {não} (nao.west);

    \node[consequencia, text width=44mm] (papeis) at (9.4,-3.5)
        {papéis: LLM Inicial pela razão acurácia/custo; LLM Avançado se
         superior com significância (senão a Fase~3 é eliminada)};
    \draw[seta] (nao) -- (papeis);

    \node[rot, text width=50mm, anchor=north] (ramo) at (9.4,-4.75)
        {o ramo efetivamente seguido, e a divergência que ele exige
         declarar, são reportados com os resultados (Seção~\ref{sec:res-gate})};
    \draw[seta, dashed] (ramo.north) -- (papeis.south);

    % ---------------- faixa inferior: a RÉGUA da hipótese ------------------
    \node[rot, anchor=west] (tit2) at (-1.9,-5.9)
        {\textbf{Critério da hipótese} (Seção~\ref{sec:intro-hipotese})};

    \node[medida] (regua) at (0,-7.4)
        {régua $D$: o mesmo\\classificador com o\\\textit{pool} de referência\\inteiramente rotulado};
    \node[decisao] (crit) at (4.6,-7.4)
        {$\mathrm{acc} \ge 0{,}95 \times \mathrm{acc}(D)$\\com no máximo\\34.724 rótulos?};
    \node[caixa] (veredito) at (9.4,-7.4)
        {veredito da hipótese\\(julgado nos resultados,\\Seção~\ref{sec:res-e3p})};

    \draw[seta] (regua) -- (crit);
    \draw[seta] (crit) -- (veredito);
    \node[rot, below=2.5mm of crit, text width=62mm]
        {teto de 34.724 rótulos $= 15\%$ da população deduplicada de 231.490;
         a acurácia é a métrica do critério (pré-registro); o Macro F1
         acompanha como análise de robustez};

    % ---------------- amarração entre as duas réguas -----------------------
    % conector vertical reto entre as duas caixas de decisão, com o rótulo
    % ancorado à esquerda da linha (nada cruza título nem seta).
    % Fica por último porque referencia nós das duas faixas.
    \draw[seta, dashed] (gate.south) -- (crit.north);
    \node[rot, text width=40mm, anchor=east] at (4.3,-3.3)
        {a mesma razão de $0{,}95$ nas duas réguas:
         $85\% \approx 0{,}95 \times 89{,}56\%$ do melhor
         baseline supervisionado leve};
\end{tikzpicture}
\end{document}
%   (no modo standalone, troque as \ref{...} por texto fixo ou defina-as)
% Requer apenas as bibliotecas já carregadas em packages.tex
% (arrows.meta, positioning). Desenhado para ser legível em preto e branco:
% decisão = borda grossa + cinza escuro; medição = cinza claro;
% consequência = cinza médio; condicional/remissão = tracejado.
% FREEZE respeitado: todos os números são os já fixados no texto
% (85%; 0,95 x 89,56%; 0,95 x acc(D); 34.724; 231.490; 15%).
% Compilado e inspecionado visualmente (pdflatex + render PNG) — ver
% NOTA-esquemas.md, iterações 4-5 do loop.
% ============================================================================
\begin{tikzpicture}[
    x=1cm, y=1cm,
    caixa/.style={draw, rounded corners=2pt, align=center, text width=34mm,
                  minimum height=11mm, inner sep=3pt, font=\footnotesize},
    decisao/.style={caixa, thick, fill=gray!25},
    medida/.style={caixa, fill=gray!8},
    consequencia/.style={caixa, fill=gray!14},
    seta/.style={-{Stealth[length=2.4mm]}, thick},
    rot/.style={font=\scriptsize, align=center, inner sep=2pt}
]
    % ---------------- faixa superior: o GATE do oráculo -------------------
    \node[rot, anchor=west] (tit1) at (-1.9,1.7)
        {\textbf{Gate do oráculo} (fixado de antemão, Seção~\ref{sec:metodo-oraculo-decisao})};

    \node[medida]  (e0)   at (0,0)
        {avaliação dos oráculos\\LLM (S-rand $n=1.000$;\\S-strat $n\approx1.863$)};
    \node[decisao] (gate) at (4.6,0)
        {algum oráculo com\\acurácia $\ge 85\%$\\na S-rand?};
    \node[caixa]   (sim)  at (9.4,1.3)
        {oráculo \textbf{aprovado}:\\papéis pelo critério\\original};
    \node[consequencia] (nao) at (9.4,-1.3)
        {robustez ao ruído\\\textbf{obrigatória}; oráculo ruidoso\\vira cenário central};

    \draw[seta] (e0) -- (gate);
    % âncoras deixam o rótulo inteiro do lado livre da diagonal
    % (a seta cortava o "m"/"o" — achado da inspeção visual)
    \draw[seta] (gate.east) -- node[rot, anchor=south east, pos=0.55] {sim} (sim.west);
    \draw[seta] (gate.east) -- node[rot, anchor=north east, pos=0.55] {não} (nao.west);

    \node[consequencia, text width=44mm] (papeis) at (9.4,-3.5)
        {papéis: LLM Inicial pela razão acurácia/custo; LLM Avançado se
         superior com significância (senão a Fase~3 é eliminada)};
    \draw[seta] (nao) -- (papeis);

    \node[rot, text width=50mm, anchor=north] (ramo) at (9.4,-4.75)
        {o ramo efetivamente seguido, e a divergência que ele exige
         declarar, são reportados com os resultados (Seção~\ref{sec:res-gate})};
    \draw[seta, dashed] (ramo.north) -- (papeis.south);

    % ---------------- faixa inferior: a RÉGUA da hipótese ------------------
    \node[rot, anchor=west] (tit2) at (-1.9,-5.9)
        {\textbf{Critério da hipótese} (Seção~\ref{sec:intro-hipotese})};

    \node[medida] (regua) at (0,-7.4)
        {régua $D$: o mesmo\\classificador com o\\\textit{pool} de referência\\inteiramente rotulado};
    \node[decisao] (crit) at (4.6,-7.4)
        {$\mathrm{acc} \ge 0{,}95 \times \mathrm{acc}(D)$\\com no máximo\\34.724 rótulos?};
    \node[caixa] (veredito) at (9.4,-7.4)
        {veredito da hipótese\\(julgado nos resultados,\\Seção~\ref{sec:res-e3p})};

    \draw[seta] (regua) -- (crit);
    \draw[seta] (crit) -- (veredito);
    \node[rot, below=2.5mm of crit, text width=62mm]
        {teto de 34.724 rótulos $= 15\%$ da população deduplicada de 231.490;
         a acurácia é a métrica do critério (pré-registro); o Macro F1
         acompanha como análise de robustez};

    % ---------------- amarração entre as duas réguas -----------------------
    % conector vertical reto entre as duas caixas de decisão, com o rótulo
    % ancorado à esquerda da linha (nada cruza título nem seta).
    % Fica por último porque referencia nós das duas faixas.
    \draw[seta, dashed] (gate.south) -- (crit.north);
    \node[rot, text width=40mm, anchor=east] at (4.3,-3.3)
        {a mesma razão de $0{,}95$ nas duas réguas:
         $85\% \approx 0{,}95 \times 89{,}56\%$ do melhor
         baseline supervisionado leve};
\end{tikzpicture}
\end{document}
%   (no modo standalone, troque as \ref{...} por texto fixo ou defina-as)
% Requer apenas as bibliotecas já carregadas em packages.tex
% (arrows.meta, positioning). Desenhado para ser legível em preto e branco:
% decisão = borda grossa + cinza escuro; medição = cinza claro;
% consequência = cinza médio; condicional/remissão = tracejado.
% FREEZE respeitado: todos os números são os já fixados no texto
% (85%; 0,95 x 89,56%; 0,95 x acc(D); 34.724; 231.490; 15%).
% Compilado e inspecionado visualmente (pdflatex + render PNG) — ver
% NOTA-esquemas.md, iterações 4-5 do loop.
% ============================================================================
\begin{tikzpicture}[
    x=1cm, y=1cm,
    caixa/.style={draw, rounded corners=2pt, align=center, text width=34mm,
                  minimum height=11mm, inner sep=3pt, font=\footnotesize},
    decisao/.style={caixa, thick, fill=gray!25},
    medida/.style={caixa, fill=gray!8},
    consequencia/.style={caixa, fill=gray!14},
    seta/.style={-{Stealth[length=2.4mm]}, thick},
    rot/.style={font=\scriptsize, align=center, inner sep=2pt}
]
    % ---------------- faixa superior: o GATE do oráculo -------------------
    \node[rot, anchor=west] (tit1) at (-1.9,1.7)
        {\textbf{Gate do oráculo} (fixado de antemão, Seção~\ref{sec:metodo-oraculo-decisao})};

    \node[medida]  (e0)   at (0,0)
        {avaliação dos oráculos\\LLM (S-rand $n=1.000$;\\S-strat $n\approx1.863$)};
    \node[decisao] (gate) at (4.6,0)
        {algum oráculo com\\acurácia $\ge 85\%$\\na S-rand?};
    \node[caixa]   (sim)  at (9.4,1.3)
        {oráculo \textbf{aprovado}:\\papéis pelo critério\\original};
    \node[consequencia] (nao) at (9.4,-1.3)
        {robustez ao ruído\\\textbf{obrigatória}; oráculo ruidoso\\vira cenário central};

    \draw[seta] (e0) -- (gate);
    % âncoras deixam o rótulo inteiro do lado livre da diagonal
    % (a seta cortava o "m"/"o" — achado da inspeção visual)
    \draw[seta] (gate.east) -- node[rot, anchor=south east, pos=0.55] {sim} (sim.west);
    \draw[seta] (gate.east) -- node[rot, anchor=north east, pos=0.55] {não} (nao.west);

    \node[consequencia, text width=44mm] (papeis) at (9.4,-3.5)
        {papéis: LLM Inicial pela razão acurácia/custo; LLM Avançado se
         superior com significância (senão a Fase~3 é eliminada)};
    \draw[seta] (nao) -- (papeis);

    \node[rot, text width=50mm, anchor=north] (ramo) at (9.4,-4.75)
        {o ramo efetivamente seguido, e a divergência que ele exige
         declarar, são reportados com os resultados (Seção~\ref{sec:res-gate})};
    \draw[seta, dashed] (ramo.north) -- (papeis.south);

    % ---------------- faixa inferior: a RÉGUA da hipótese ------------------
    \node[rot, anchor=west] (tit2) at (-1.9,-5.9)
        {\textbf{Critério da hipótese} (Seção~\ref{sec:intro-hipotese})};

    \node[medida] (regua) at (0,-7.4)
        {régua $D$: o mesmo\\classificador com o\\\textit{pool} de referência\\inteiramente rotulado};
    \node[decisao] (crit) at (4.6,-7.4)
        {$\mathrm{acc} \ge 0{,}95 \times \mathrm{acc}(D)$\\com no máximo\\34.724 rótulos?};
    \node[caixa] (veredito) at (9.4,-7.4)
        {veredito da hipótese\\(julgado nos resultados,\\Seção~\ref{sec:res-e3p})};

    \draw[seta] (regua) -- (crit);
    \draw[seta] (crit) -- (veredito);
    \node[rot, below=2.5mm of crit, text width=62mm]
        {teto de 34.724 rótulos $= 15\%$ da população deduplicada de 231.490;
         a acurácia é a métrica do critério (pré-registro); o Macro F1
         acompanha como análise de robustez};

    % ---------------- amarração entre as duas réguas -----------------------
    % conector vertical reto entre as duas caixas de decisão, com o rótulo
    % ancorado à esquerda da linha (nada cruza título nem seta).
    % Fica por último porque referencia nós das duas faixas.
    \draw[seta, dashed] (gate.south) -- (crit.north);
    \node[rot, text width=40mm, anchor=east] at (4.3,-3.3)
        {a mesma razão de $0{,}95$ nas duas réguas:
         $85\% \approx 0{,}95 \times 89{,}56\%$ do melhor
         baseline supervisionado leve};
\end{tikzpicture}

% Uso standalone (pré-visualização):
%   \documentclass[tikz,border=6pt]{standalone}
%   \usetikzlibrary{arrows.meta, positioning}
%   \begin{document}% ============================================================================
% ESQUEMA PROPOSTO (banca, sugestão — não mergeado): o gate do oráculo e a
% régua da hipótese como fluxo de decisão. Complementa a Seção
% sec:metodo-oraculo-decisao e a Seção sec:intro-hipotese.
%
% Uso no Cap. 3 (dentro de um ambiente figure):
%   % ============================================================================
% ESQUEMA PROPOSTO (banca, sugestão — não mergeado): o gate do oráculo e a
% régua da hipótese como fluxo de decisão. Complementa a Seção
% sec:metodo-oraculo-decisao e a Seção sec:intro-hipotese.
%
% Uso no Cap. 3 (dentro de um ambiente figure):
%   % ============================================================================
% ESQUEMA PROPOSTO (banca, sugestão — não mergeado): o gate do oráculo e a
% régua da hipótese como fluxo de decisão. Complementa a Seção
% sec:metodo-oraculo-decisao e a Seção sec:intro-hipotese.
%
% Uso no Cap. 3 (dentro de um ambiente figure):
%   \input{3-metodo/esquemas-propostos/esq-gate-e-regua}
% Uso standalone (pré-visualização):
%   \documentclass[tikz,border=6pt]{standalone}
%   \usetikzlibrary{arrows.meta, positioning}
%   \begin{document}\input{esq-gate-e-regua.tex}\end{document}
%   (no modo standalone, troque as \ref{...} por texto fixo ou defina-as)
% Requer apenas as bibliotecas já carregadas em packages.tex
% (arrows.meta, positioning). Desenhado para ser legível em preto e branco:
% decisão = borda grossa + cinza escuro; medição = cinza claro;
% consequência = cinza médio; condicional/remissão = tracejado.
% FREEZE respeitado: todos os números são os já fixados no texto
% (85%; 0,95 x 89,56%; 0,95 x acc(D); 34.724; 231.490; 15%).
% Compilado e inspecionado visualmente (pdflatex + render PNG) — ver
% NOTA-esquemas.md, iterações 4-5 do loop.
% ============================================================================
\begin{tikzpicture}[
    x=1cm, y=1cm,
    caixa/.style={draw, rounded corners=2pt, align=center, text width=34mm,
                  minimum height=11mm, inner sep=3pt, font=\footnotesize},
    decisao/.style={caixa, thick, fill=gray!25},
    medida/.style={caixa, fill=gray!8},
    consequencia/.style={caixa, fill=gray!14},
    seta/.style={-{Stealth[length=2.4mm]}, thick},
    rot/.style={font=\scriptsize, align=center, inner sep=2pt}
]
    % ---------------- faixa superior: o GATE do oráculo -------------------
    \node[rot, anchor=west] (tit1) at (-1.9,1.7)
        {\textbf{Gate do oráculo} (fixado de antemão, Seção~\ref{sec:metodo-oraculo-decisao})};

    \node[medida]  (e0)   at (0,0)
        {avaliação dos oráculos\\LLM (S-rand $n=1.000$;\\S-strat $n\approx1.863$)};
    \node[decisao] (gate) at (4.6,0)
        {algum oráculo com\\acurácia $\ge 85\%$\\na S-rand?};
    \node[caixa]   (sim)  at (9.4,1.3)
        {oráculo \textbf{aprovado}:\\papéis pelo critério\\original};
    \node[consequencia] (nao) at (9.4,-1.3)
        {robustez ao ruído\\\textbf{obrigatória}; oráculo ruidoso\\vira cenário central};

    \draw[seta] (e0) -- (gate);
    % âncoras deixam o rótulo inteiro do lado livre da diagonal
    % (a seta cortava o "m"/"o" — achado da inspeção visual)
    \draw[seta] (gate.east) -- node[rot, anchor=south east, pos=0.55] {sim} (sim.west);
    \draw[seta] (gate.east) -- node[rot, anchor=north east, pos=0.55] {não} (nao.west);

    \node[consequencia, text width=44mm] (papeis) at (9.4,-3.5)
        {papéis: LLM Inicial pela razão acurácia/custo; LLM Avançado se
         superior com significância (senão a Fase~3 é eliminada)};
    \draw[seta] (nao) -- (papeis);

    \node[rot, text width=50mm, anchor=north] (ramo) at (9.4,-4.75)
        {o ramo efetivamente seguido, e a divergência que ele exige
         declarar, são reportados com os resultados (Seção~\ref{sec:res-gate})};
    \draw[seta, dashed] (ramo.north) -- (papeis.south);

    % ---------------- faixa inferior: a RÉGUA da hipótese ------------------
    \node[rot, anchor=west] (tit2) at (-1.9,-5.9)
        {\textbf{Critério da hipótese} (Seção~\ref{sec:intro-hipotese})};

    \node[medida] (regua) at (0,-7.4)
        {régua $D$: o mesmo\\classificador com o\\\textit{pool} de referência\\inteiramente rotulado};
    \node[decisao] (crit) at (4.6,-7.4)
        {$\mathrm{acc} \ge 0{,}95 \times \mathrm{acc}(D)$\\com no máximo\\34.724 rótulos?};
    \node[caixa] (veredito) at (9.4,-7.4)
        {veredito da hipótese\\(julgado nos resultados,\\Seção~\ref{sec:res-e3p})};

    \draw[seta] (regua) -- (crit);
    \draw[seta] (crit) -- (veredito);
    \node[rot, below=2.5mm of crit, text width=62mm]
        {teto de 34.724 rótulos $= 15\%$ da população deduplicada de 231.490;
         a acurácia é a métrica do critério (pré-registro); o Macro F1
         acompanha como análise de robustez};

    % ---------------- amarração entre as duas réguas -----------------------
    % conector vertical reto entre as duas caixas de decisão, com o rótulo
    % ancorado à esquerda da linha (nada cruza título nem seta).
    % Fica por último porque referencia nós das duas faixas.
    \draw[seta, dashed] (gate.south) -- (crit.north);
    \node[rot, text width=40mm, anchor=east] at (4.3,-3.3)
        {a mesma razão de $0{,}95$ nas duas réguas:
         $85\% \approx 0{,}95 \times 89{,}56\%$ do melhor
         baseline supervisionado leve};
\end{tikzpicture}

% Uso standalone (pré-visualização):
%   \documentclass[tikz,border=6pt]{standalone}
%   \usetikzlibrary{arrows.meta, positioning}
%   \begin{document}% ============================================================================
% ESQUEMA PROPOSTO (banca, sugestão — não mergeado): o gate do oráculo e a
% régua da hipótese como fluxo de decisão. Complementa a Seção
% sec:metodo-oraculo-decisao e a Seção sec:intro-hipotese.
%
% Uso no Cap. 3 (dentro de um ambiente figure):
%   \input{3-metodo/esquemas-propostos/esq-gate-e-regua}
% Uso standalone (pré-visualização):
%   \documentclass[tikz,border=6pt]{standalone}
%   \usetikzlibrary{arrows.meta, positioning}
%   \begin{document}\input{esq-gate-e-regua.tex}\end{document}
%   (no modo standalone, troque as \ref{...} por texto fixo ou defina-as)
% Requer apenas as bibliotecas já carregadas em packages.tex
% (arrows.meta, positioning). Desenhado para ser legível em preto e branco:
% decisão = borda grossa + cinza escuro; medição = cinza claro;
% consequência = cinza médio; condicional/remissão = tracejado.
% FREEZE respeitado: todos os números são os já fixados no texto
% (85%; 0,95 x 89,56%; 0,95 x acc(D); 34.724; 231.490; 15%).
% Compilado e inspecionado visualmente (pdflatex + render PNG) — ver
% NOTA-esquemas.md, iterações 4-5 do loop.
% ============================================================================
\begin{tikzpicture}[
    x=1cm, y=1cm,
    caixa/.style={draw, rounded corners=2pt, align=center, text width=34mm,
                  minimum height=11mm, inner sep=3pt, font=\footnotesize},
    decisao/.style={caixa, thick, fill=gray!25},
    medida/.style={caixa, fill=gray!8},
    consequencia/.style={caixa, fill=gray!14},
    seta/.style={-{Stealth[length=2.4mm]}, thick},
    rot/.style={font=\scriptsize, align=center, inner sep=2pt}
]
    % ---------------- faixa superior: o GATE do oráculo -------------------
    \node[rot, anchor=west] (tit1) at (-1.9,1.7)
        {\textbf{Gate do oráculo} (fixado de antemão, Seção~\ref{sec:metodo-oraculo-decisao})};

    \node[medida]  (e0)   at (0,0)
        {avaliação dos oráculos\\LLM (S-rand $n=1.000$;\\S-strat $n\approx1.863$)};
    \node[decisao] (gate) at (4.6,0)
        {algum oráculo com\\acurácia $\ge 85\%$\\na S-rand?};
    \node[caixa]   (sim)  at (9.4,1.3)
        {oráculo \textbf{aprovado}:\\papéis pelo critério\\original};
    \node[consequencia] (nao) at (9.4,-1.3)
        {robustez ao ruído\\\textbf{obrigatória}; oráculo ruidoso\\vira cenário central};

    \draw[seta] (e0) -- (gate);
    % âncoras deixam o rótulo inteiro do lado livre da diagonal
    % (a seta cortava o "m"/"o" — achado da inspeção visual)
    \draw[seta] (gate.east) -- node[rot, anchor=south east, pos=0.55] {sim} (sim.west);
    \draw[seta] (gate.east) -- node[rot, anchor=north east, pos=0.55] {não} (nao.west);

    \node[consequencia, text width=44mm] (papeis) at (9.4,-3.5)
        {papéis: LLM Inicial pela razão acurácia/custo; LLM Avançado se
         superior com significância (senão a Fase~3 é eliminada)};
    \draw[seta] (nao) -- (papeis);

    \node[rot, text width=50mm, anchor=north] (ramo) at (9.4,-4.75)
        {o ramo efetivamente seguido, e a divergência que ele exige
         declarar, são reportados com os resultados (Seção~\ref{sec:res-gate})};
    \draw[seta, dashed] (ramo.north) -- (papeis.south);

    % ---------------- faixa inferior: a RÉGUA da hipótese ------------------
    \node[rot, anchor=west] (tit2) at (-1.9,-5.9)
        {\textbf{Critério da hipótese} (Seção~\ref{sec:intro-hipotese})};

    \node[medida] (regua) at (0,-7.4)
        {régua $D$: o mesmo\\classificador com o\\\textit{pool} de referência\\inteiramente rotulado};
    \node[decisao] (crit) at (4.6,-7.4)
        {$\mathrm{acc} \ge 0{,}95 \times \mathrm{acc}(D)$\\com no máximo\\34.724 rótulos?};
    \node[caixa] (veredito) at (9.4,-7.4)
        {veredito da hipótese\\(julgado nos resultados,\\Seção~\ref{sec:res-e3p})};

    \draw[seta] (regua) -- (crit);
    \draw[seta] (crit) -- (veredito);
    \node[rot, below=2.5mm of crit, text width=62mm]
        {teto de 34.724 rótulos $= 15\%$ da população deduplicada de 231.490;
         a acurácia é a métrica do critério (pré-registro); o Macro F1
         acompanha como análise de robustez};

    % ---------------- amarração entre as duas réguas -----------------------
    % conector vertical reto entre as duas caixas de decisão, com o rótulo
    % ancorado à esquerda da linha (nada cruza título nem seta).
    % Fica por último porque referencia nós das duas faixas.
    \draw[seta, dashed] (gate.south) -- (crit.north);
    \node[rot, text width=40mm, anchor=east] at (4.3,-3.3)
        {a mesma razão de $0{,}95$ nas duas réguas:
         $85\% \approx 0{,}95 \times 89{,}56\%$ do melhor
         baseline supervisionado leve};
\end{tikzpicture}
\end{document}
%   (no modo standalone, troque as \ref{...} por texto fixo ou defina-as)
% Requer apenas as bibliotecas já carregadas em packages.tex
% (arrows.meta, positioning). Desenhado para ser legível em preto e branco:
% decisão = borda grossa + cinza escuro; medição = cinza claro;
% consequência = cinza médio; condicional/remissão = tracejado.
% FREEZE respeitado: todos os números são os já fixados no texto
% (85%; 0,95 x 89,56%; 0,95 x acc(D); 34.724; 231.490; 15%).
% Compilado e inspecionado visualmente (pdflatex + render PNG) — ver
% NOTA-esquemas.md, iterações 4-5 do loop.
% ============================================================================
\begin{tikzpicture}[
    x=1cm, y=1cm,
    caixa/.style={draw, rounded corners=2pt, align=center, text width=34mm,
                  minimum height=11mm, inner sep=3pt, font=\footnotesize},
    decisao/.style={caixa, thick, fill=gray!25},
    medida/.style={caixa, fill=gray!8},
    consequencia/.style={caixa, fill=gray!14},
    seta/.style={-{Stealth[length=2.4mm]}, thick},
    rot/.style={font=\scriptsize, align=center, inner sep=2pt}
]
    % ---------------- faixa superior: o GATE do oráculo -------------------
    \node[rot, anchor=west] (tit1) at (-1.9,1.7)
        {\textbf{Gate do oráculo} (fixado de antemão, Seção~\ref{sec:metodo-oraculo-decisao})};

    \node[medida]  (e0)   at (0,0)
        {avaliação dos oráculos\\LLM (S-rand $n=1.000$;\\S-strat $n\approx1.863$)};
    \node[decisao] (gate) at (4.6,0)
        {algum oráculo com\\acurácia $\ge 85\%$\\na S-rand?};
    \node[caixa]   (sim)  at (9.4,1.3)
        {oráculo \textbf{aprovado}:\\papéis pelo critério\\original};
    \node[consequencia] (nao) at (9.4,-1.3)
        {robustez ao ruído\\\textbf{obrigatória}; oráculo ruidoso\\vira cenário central};

    \draw[seta] (e0) -- (gate);
    % âncoras deixam o rótulo inteiro do lado livre da diagonal
    % (a seta cortava o "m"/"o" — achado da inspeção visual)
    \draw[seta] (gate.east) -- node[rot, anchor=south east, pos=0.55] {sim} (sim.west);
    \draw[seta] (gate.east) -- node[rot, anchor=north east, pos=0.55] {não} (nao.west);

    \node[consequencia, text width=44mm] (papeis) at (9.4,-3.5)
        {papéis: LLM Inicial pela razão acurácia/custo; LLM Avançado se
         superior com significância (senão a Fase~3 é eliminada)};
    \draw[seta] (nao) -- (papeis);

    \node[rot, text width=50mm, anchor=north] (ramo) at (9.4,-4.75)
        {o ramo efetivamente seguido, e a divergência que ele exige
         declarar, são reportados com os resultados (Seção~\ref{sec:res-gate})};
    \draw[seta, dashed] (ramo.north) -- (papeis.south);

    % ---------------- faixa inferior: a RÉGUA da hipótese ------------------
    \node[rot, anchor=west] (tit2) at (-1.9,-5.9)
        {\textbf{Critério da hipótese} (Seção~\ref{sec:intro-hipotese})};

    \node[medida] (regua) at (0,-7.4)
        {régua $D$: o mesmo\\classificador com o\\\textit{pool} de referência\\inteiramente rotulado};
    \node[decisao] (crit) at (4.6,-7.4)
        {$\mathrm{acc} \ge 0{,}95 \times \mathrm{acc}(D)$\\com no máximo\\34.724 rótulos?};
    \node[caixa] (veredito) at (9.4,-7.4)
        {veredito da hipótese\\(julgado nos resultados,\\Seção~\ref{sec:res-e3p})};

    \draw[seta] (regua) -- (crit);
    \draw[seta] (crit) -- (veredito);
    \node[rot, below=2.5mm of crit, text width=62mm]
        {teto de 34.724 rótulos $= 15\%$ da população deduplicada de 231.490;
         a acurácia é a métrica do critério (pré-registro); o Macro F1
         acompanha como análise de robustez};

    % ---------------- amarração entre as duas réguas -----------------------
    % conector vertical reto entre as duas caixas de decisão, com o rótulo
    % ancorado à esquerda da linha (nada cruza título nem seta).
    % Fica por último porque referencia nós das duas faixas.
    \draw[seta, dashed] (gate.south) -- (crit.north);
    \node[rot, text width=40mm, anchor=east] at (4.3,-3.3)
        {a mesma razão de $0{,}95$ nas duas réguas:
         $85\% \approx 0{,}95 \times 89{,}56\%$ do melhor
         baseline supervisionado leve};
\end{tikzpicture}

% Uso standalone (pré-visualização):
%   \documentclass[tikz,border=6pt]{standalone}
%   \usetikzlibrary{arrows.meta, positioning}
%   \begin{document}% ============================================================================
% ESQUEMA PROPOSTO (banca, sugestão — não mergeado): o gate do oráculo e a
% régua da hipótese como fluxo de decisão. Complementa a Seção
% sec:metodo-oraculo-decisao e a Seção sec:intro-hipotese.
%
% Uso no Cap. 3 (dentro de um ambiente figure):
%   % ============================================================================
% ESQUEMA PROPOSTO (banca, sugestão — não mergeado): o gate do oráculo e a
% régua da hipótese como fluxo de decisão. Complementa a Seção
% sec:metodo-oraculo-decisao e a Seção sec:intro-hipotese.
%
% Uso no Cap. 3 (dentro de um ambiente figure):
%   \input{3-metodo/esquemas-propostos/esq-gate-e-regua}
% Uso standalone (pré-visualização):
%   \documentclass[tikz,border=6pt]{standalone}
%   \usetikzlibrary{arrows.meta, positioning}
%   \begin{document}\input{esq-gate-e-regua.tex}\end{document}
%   (no modo standalone, troque as \ref{...} por texto fixo ou defina-as)
% Requer apenas as bibliotecas já carregadas em packages.tex
% (arrows.meta, positioning). Desenhado para ser legível em preto e branco:
% decisão = borda grossa + cinza escuro; medição = cinza claro;
% consequência = cinza médio; condicional/remissão = tracejado.
% FREEZE respeitado: todos os números são os já fixados no texto
% (85%; 0,95 x 89,56%; 0,95 x acc(D); 34.724; 231.490; 15%).
% Compilado e inspecionado visualmente (pdflatex + render PNG) — ver
% NOTA-esquemas.md, iterações 4-5 do loop.
% ============================================================================
\begin{tikzpicture}[
    x=1cm, y=1cm,
    caixa/.style={draw, rounded corners=2pt, align=center, text width=34mm,
                  minimum height=11mm, inner sep=3pt, font=\footnotesize},
    decisao/.style={caixa, thick, fill=gray!25},
    medida/.style={caixa, fill=gray!8},
    consequencia/.style={caixa, fill=gray!14},
    seta/.style={-{Stealth[length=2.4mm]}, thick},
    rot/.style={font=\scriptsize, align=center, inner sep=2pt}
]
    % ---------------- faixa superior: o GATE do oráculo -------------------
    \node[rot, anchor=west] (tit1) at (-1.9,1.7)
        {\textbf{Gate do oráculo} (fixado de antemão, Seção~\ref{sec:metodo-oraculo-decisao})};

    \node[medida]  (e0)   at (0,0)
        {avaliação dos oráculos\\LLM (S-rand $n=1.000$;\\S-strat $n\approx1.863$)};
    \node[decisao] (gate) at (4.6,0)
        {algum oráculo com\\acurácia $\ge 85\%$\\na S-rand?};
    \node[caixa]   (sim)  at (9.4,1.3)
        {oráculo \textbf{aprovado}:\\papéis pelo critério\\original};
    \node[consequencia] (nao) at (9.4,-1.3)
        {robustez ao ruído\\\textbf{obrigatória}; oráculo ruidoso\\vira cenário central};

    \draw[seta] (e0) -- (gate);
    % âncoras deixam o rótulo inteiro do lado livre da diagonal
    % (a seta cortava o "m"/"o" — achado da inspeção visual)
    \draw[seta] (gate.east) -- node[rot, anchor=south east, pos=0.55] {sim} (sim.west);
    \draw[seta] (gate.east) -- node[rot, anchor=north east, pos=0.55] {não} (nao.west);

    \node[consequencia, text width=44mm] (papeis) at (9.4,-3.5)
        {papéis: LLM Inicial pela razão acurácia/custo; LLM Avançado se
         superior com significância (senão a Fase~3 é eliminada)};
    \draw[seta] (nao) -- (papeis);

    \node[rot, text width=50mm, anchor=north] (ramo) at (9.4,-4.75)
        {o ramo efetivamente seguido, e a divergência que ele exige
         declarar, são reportados com os resultados (Seção~\ref{sec:res-gate})};
    \draw[seta, dashed] (ramo.north) -- (papeis.south);

    % ---------------- faixa inferior: a RÉGUA da hipótese ------------------
    \node[rot, anchor=west] (tit2) at (-1.9,-5.9)
        {\textbf{Critério da hipótese} (Seção~\ref{sec:intro-hipotese})};

    \node[medida] (regua) at (0,-7.4)
        {régua $D$: o mesmo\\classificador com o\\\textit{pool} de referência\\inteiramente rotulado};
    \node[decisao] (crit) at (4.6,-7.4)
        {$\mathrm{acc} \ge 0{,}95 \times \mathrm{acc}(D)$\\com no máximo\\34.724 rótulos?};
    \node[caixa] (veredito) at (9.4,-7.4)
        {veredito da hipótese\\(julgado nos resultados,\\Seção~\ref{sec:res-e3p})};

    \draw[seta] (regua) -- (crit);
    \draw[seta] (crit) -- (veredito);
    \node[rot, below=2.5mm of crit, text width=62mm]
        {teto de 34.724 rótulos $= 15\%$ da população deduplicada de 231.490;
         a acurácia é a métrica do critério (pré-registro); o Macro F1
         acompanha como análise de robustez};

    % ---------------- amarração entre as duas réguas -----------------------
    % conector vertical reto entre as duas caixas de decisão, com o rótulo
    % ancorado à esquerda da linha (nada cruza título nem seta).
    % Fica por último porque referencia nós das duas faixas.
    \draw[seta, dashed] (gate.south) -- (crit.north);
    \node[rot, text width=40mm, anchor=east] at (4.3,-3.3)
        {a mesma razão de $0{,}95$ nas duas réguas:
         $85\% \approx 0{,}95 \times 89{,}56\%$ do melhor
         baseline supervisionado leve};
\end{tikzpicture}

% Uso standalone (pré-visualização):
%   \documentclass[tikz,border=6pt]{standalone}
%   \usetikzlibrary{arrows.meta, positioning}
%   \begin{document}% ============================================================================
% ESQUEMA PROPOSTO (banca, sugestão — não mergeado): o gate do oráculo e a
% régua da hipótese como fluxo de decisão. Complementa a Seção
% sec:metodo-oraculo-decisao e a Seção sec:intro-hipotese.
%
% Uso no Cap. 3 (dentro de um ambiente figure):
%   \input{3-metodo/esquemas-propostos/esq-gate-e-regua}
% Uso standalone (pré-visualização):
%   \documentclass[tikz,border=6pt]{standalone}
%   \usetikzlibrary{arrows.meta, positioning}
%   \begin{document}\input{esq-gate-e-regua.tex}\end{document}
%   (no modo standalone, troque as \ref{...} por texto fixo ou defina-as)
% Requer apenas as bibliotecas já carregadas em packages.tex
% (arrows.meta, positioning). Desenhado para ser legível em preto e branco:
% decisão = borda grossa + cinza escuro; medição = cinza claro;
% consequência = cinza médio; condicional/remissão = tracejado.
% FREEZE respeitado: todos os números são os já fixados no texto
% (85%; 0,95 x 89,56%; 0,95 x acc(D); 34.724; 231.490; 15%).
% Compilado e inspecionado visualmente (pdflatex + render PNG) — ver
% NOTA-esquemas.md, iterações 4-5 do loop.
% ============================================================================
\begin{tikzpicture}[
    x=1cm, y=1cm,
    caixa/.style={draw, rounded corners=2pt, align=center, text width=34mm,
                  minimum height=11mm, inner sep=3pt, font=\footnotesize},
    decisao/.style={caixa, thick, fill=gray!25},
    medida/.style={caixa, fill=gray!8},
    consequencia/.style={caixa, fill=gray!14},
    seta/.style={-{Stealth[length=2.4mm]}, thick},
    rot/.style={font=\scriptsize, align=center, inner sep=2pt}
]
    % ---------------- faixa superior: o GATE do oráculo -------------------
    \node[rot, anchor=west] (tit1) at (-1.9,1.7)
        {\textbf{Gate do oráculo} (fixado de antemão, Seção~\ref{sec:metodo-oraculo-decisao})};

    \node[medida]  (e0)   at (0,0)
        {avaliação dos oráculos\\LLM (S-rand $n=1.000$;\\S-strat $n\approx1.863$)};
    \node[decisao] (gate) at (4.6,0)
        {algum oráculo com\\acurácia $\ge 85\%$\\na S-rand?};
    \node[caixa]   (sim)  at (9.4,1.3)
        {oráculo \textbf{aprovado}:\\papéis pelo critério\\original};
    \node[consequencia] (nao) at (9.4,-1.3)
        {robustez ao ruído\\\textbf{obrigatória}; oráculo ruidoso\\vira cenário central};

    \draw[seta] (e0) -- (gate);
    % âncoras deixam o rótulo inteiro do lado livre da diagonal
    % (a seta cortava o "m"/"o" — achado da inspeção visual)
    \draw[seta] (gate.east) -- node[rot, anchor=south east, pos=0.55] {sim} (sim.west);
    \draw[seta] (gate.east) -- node[rot, anchor=north east, pos=0.55] {não} (nao.west);

    \node[consequencia, text width=44mm] (papeis) at (9.4,-3.5)
        {papéis: LLM Inicial pela razão acurácia/custo; LLM Avançado se
         superior com significância (senão a Fase~3 é eliminada)};
    \draw[seta] (nao) -- (papeis);

    \node[rot, text width=50mm, anchor=north] (ramo) at (9.4,-4.75)
        {o ramo efetivamente seguido, e a divergência que ele exige
         declarar, são reportados com os resultados (Seção~\ref{sec:res-gate})};
    \draw[seta, dashed] (ramo.north) -- (papeis.south);

    % ---------------- faixa inferior: a RÉGUA da hipótese ------------------
    \node[rot, anchor=west] (tit2) at (-1.9,-5.9)
        {\textbf{Critério da hipótese} (Seção~\ref{sec:intro-hipotese})};

    \node[medida] (regua) at (0,-7.4)
        {régua $D$: o mesmo\\classificador com o\\\textit{pool} de referência\\inteiramente rotulado};
    \node[decisao] (crit) at (4.6,-7.4)
        {$\mathrm{acc} \ge 0{,}95 \times \mathrm{acc}(D)$\\com no máximo\\34.724 rótulos?};
    \node[caixa] (veredito) at (9.4,-7.4)
        {veredito da hipótese\\(julgado nos resultados,\\Seção~\ref{sec:res-e3p})};

    \draw[seta] (regua) -- (crit);
    \draw[seta] (crit) -- (veredito);
    \node[rot, below=2.5mm of crit, text width=62mm]
        {teto de 34.724 rótulos $= 15\%$ da população deduplicada de 231.490;
         a acurácia é a métrica do critério (pré-registro); o Macro F1
         acompanha como análise de robustez};

    % ---------------- amarração entre as duas réguas -----------------------
    % conector vertical reto entre as duas caixas de decisão, com o rótulo
    % ancorado à esquerda da linha (nada cruza título nem seta).
    % Fica por último porque referencia nós das duas faixas.
    \draw[seta, dashed] (gate.south) -- (crit.north);
    \node[rot, text width=40mm, anchor=east] at (4.3,-3.3)
        {a mesma razão de $0{,}95$ nas duas réguas:
         $85\% \approx 0{,}95 \times 89{,}56\%$ do melhor
         baseline supervisionado leve};
\end{tikzpicture}
\end{document}
%   (no modo standalone, troque as \ref{...} por texto fixo ou defina-as)
% Requer apenas as bibliotecas já carregadas em packages.tex
% (arrows.meta, positioning). Desenhado para ser legível em preto e branco:
% decisão = borda grossa + cinza escuro; medição = cinza claro;
% consequência = cinza médio; condicional/remissão = tracejado.
% FREEZE respeitado: todos os números são os já fixados no texto
% (85%; 0,95 x 89,56%; 0,95 x acc(D); 34.724; 231.490; 15%).
% Compilado e inspecionado visualmente (pdflatex + render PNG) — ver
% NOTA-esquemas.md, iterações 4-5 do loop.
% ============================================================================
\begin{tikzpicture}[
    x=1cm, y=1cm,
    caixa/.style={draw, rounded corners=2pt, align=center, text width=34mm,
                  minimum height=11mm, inner sep=3pt, font=\footnotesize},
    decisao/.style={caixa, thick, fill=gray!25},
    medida/.style={caixa, fill=gray!8},
    consequencia/.style={caixa, fill=gray!14},
    seta/.style={-{Stealth[length=2.4mm]}, thick},
    rot/.style={font=\scriptsize, align=center, inner sep=2pt}
]
    % ---------------- faixa superior: o GATE do oráculo -------------------
    \node[rot, anchor=west] (tit1) at (-1.9,1.7)
        {\textbf{Gate do oráculo} (fixado de antemão, Seção~\ref{sec:metodo-oraculo-decisao})};

    \node[medida]  (e0)   at (0,0)
        {avaliação dos oráculos\\LLM (S-rand $n=1.000$;\\S-strat $n\approx1.863$)};
    \node[decisao] (gate) at (4.6,0)
        {algum oráculo com\\acurácia $\ge 85\%$\\na S-rand?};
    \node[caixa]   (sim)  at (9.4,1.3)
        {oráculo \textbf{aprovado}:\\papéis pelo critério\\original};
    \node[consequencia] (nao) at (9.4,-1.3)
        {robustez ao ruído\\\textbf{obrigatória}; oráculo ruidoso\\vira cenário central};

    \draw[seta] (e0) -- (gate);
    % âncoras deixam o rótulo inteiro do lado livre da diagonal
    % (a seta cortava o "m"/"o" — achado da inspeção visual)
    \draw[seta] (gate.east) -- node[rot, anchor=south east, pos=0.55] {sim} (sim.west);
    \draw[seta] (gate.east) -- node[rot, anchor=north east, pos=0.55] {não} (nao.west);

    \node[consequencia, text width=44mm] (papeis) at (9.4,-3.5)
        {papéis: LLM Inicial pela razão acurácia/custo; LLM Avançado se
         superior com significância (senão a Fase~3 é eliminada)};
    \draw[seta] (nao) -- (papeis);

    \node[rot, text width=50mm, anchor=north] (ramo) at (9.4,-4.75)
        {o ramo efetivamente seguido, e a divergência que ele exige
         declarar, são reportados com os resultados (Seção~\ref{sec:res-gate})};
    \draw[seta, dashed] (ramo.north) -- (papeis.south);

    % ---------------- faixa inferior: a RÉGUA da hipótese ------------------
    \node[rot, anchor=west] (tit2) at (-1.9,-5.9)
        {\textbf{Critério da hipótese} (Seção~\ref{sec:intro-hipotese})};

    \node[medida] (regua) at (0,-7.4)
        {régua $D$: o mesmo\\classificador com o\\\textit{pool} de referência\\inteiramente rotulado};
    \node[decisao] (crit) at (4.6,-7.4)
        {$\mathrm{acc} \ge 0{,}95 \times \mathrm{acc}(D)$\\com no máximo\\34.724 rótulos?};
    \node[caixa] (veredito) at (9.4,-7.4)
        {veredito da hipótese\\(julgado nos resultados,\\Seção~\ref{sec:res-e3p})};

    \draw[seta] (regua) -- (crit);
    \draw[seta] (crit) -- (veredito);
    \node[rot, below=2.5mm of crit, text width=62mm]
        {teto de 34.724 rótulos $= 15\%$ da população deduplicada de 231.490;
         a acurácia é a métrica do critério (pré-registro); o Macro F1
         acompanha como análise de robustez};

    % ---------------- amarração entre as duas réguas -----------------------
    % conector vertical reto entre as duas caixas de decisão, com o rótulo
    % ancorado à esquerda da linha (nada cruza título nem seta).
    % Fica por último porque referencia nós das duas faixas.
    \draw[seta, dashed] (gate.south) -- (crit.north);
    \node[rot, text width=40mm, anchor=east] at (4.3,-3.3)
        {a mesma razão de $0{,}95$ nas duas réguas:
         $85\% \approx 0{,}95 \times 89{,}56\%$ do melhor
         baseline supervisionado leve};
\end{tikzpicture}
\end{document}
%   (no modo standalone, troque as \ref{...} por texto fixo ou defina-as)
% Requer apenas as bibliotecas já carregadas em packages.tex
% (arrows.meta, positioning). Desenhado para ser legível em preto e branco:
% decisão = borda grossa + cinza escuro; medição = cinza claro;
% consequência = cinza médio; condicional/remissão = tracejado.
% FREEZE respeitado: todos os números são os já fixados no texto
% (85%; 0,95 x 89,56%; 0,95 x acc(D); 34.724; 231.490; 15%).
% Compilado e inspecionado visualmente (pdflatex + render PNG) — ver
% NOTA-esquemas.md, iterações 4-5 do loop.
% ============================================================================
\begin{tikzpicture}[
    x=1cm, y=1cm,
    caixa/.style={draw, rounded corners=2pt, align=center, text width=34mm,
                  minimum height=11mm, inner sep=3pt, font=\footnotesize},
    decisao/.style={caixa, thick, fill=gray!25},
    medida/.style={caixa, fill=gray!8},
    consequencia/.style={caixa, fill=gray!14},
    seta/.style={-{Stealth[length=2.4mm]}, thick},
    rot/.style={font=\scriptsize, align=center, inner sep=2pt}
]
    % ---------------- faixa superior: o GATE do oráculo -------------------
    \node[rot, anchor=west] (tit1) at (-1.9,1.7)
        {\textbf{Gate do oráculo} (fixado de antemão, Seção~\ref{sec:metodo-oraculo-decisao})};

    \node[medida]  (e0)   at (0,0)
        {avaliação dos oráculos\\LLM (S-rand $n=1.000$;\\S-strat $n\approx1.863$)};
    \node[decisao] (gate) at (4.6,0)
        {algum oráculo com\\acurácia $\ge 85\%$\\na S-rand?};
    \node[caixa]   (sim)  at (9.4,1.3)
        {oráculo \textbf{aprovado}:\\papéis pelo critério\\original};
    \node[consequencia] (nao) at (9.4,-1.3)
        {robustez ao ruído\\\textbf{obrigatória}; oráculo ruidoso\\vira cenário central};

    \draw[seta] (e0) -- (gate);
    % âncoras deixam o rótulo inteiro do lado livre da diagonal
    % (a seta cortava o "m"/"o" — achado da inspeção visual)
    \draw[seta] (gate.east) -- node[rot, anchor=south east, pos=0.55] {sim} (sim.west);
    \draw[seta] (gate.east) -- node[rot, anchor=north east, pos=0.55] {não} (nao.west);

    \node[consequencia, text width=44mm] (papeis) at (9.4,-3.5)
        {papéis: LLM Inicial pela razão acurácia/custo; LLM Avançado se
         superior com significância (senão a Fase~3 é eliminada)};
    \draw[seta] (nao) -- (papeis);

    \node[rot, text width=50mm, anchor=north] (ramo) at (9.4,-4.75)
        {o ramo efetivamente seguido, e a divergência que ele exige
         declarar, são reportados com os resultados (Seção~\ref{sec:res-gate})};
    \draw[seta, dashed] (ramo.north) -- (papeis.south);

    % ---------------- faixa inferior: a RÉGUA da hipótese ------------------
    \node[rot, anchor=west] (tit2) at (-1.9,-5.9)
        {\textbf{Critério da hipótese} (Seção~\ref{sec:intro-hipotese})};

    \node[medida] (regua) at (0,-7.4)
        {régua $D$: o mesmo\\classificador com o\\\textit{pool} de referência\\inteiramente rotulado};
    \node[decisao] (crit) at (4.6,-7.4)
        {$\mathrm{acc} \ge 0{,}95 \times \mathrm{acc}(D)$\\com no máximo\\34.724 rótulos?};
    \node[caixa] (veredito) at (9.4,-7.4)
        {veredito da hipótese\\(julgado nos resultados,\\Seção~\ref{sec:res-e3p})};

    \draw[seta] (regua) -- (crit);
    \draw[seta] (crit) -- (veredito);
    \node[rot, below=2.5mm of crit, text width=62mm]
        {teto de 34.724 rótulos $= 15\%$ da população deduplicada de 231.490;
         a acurácia é a métrica do critério (pré-registro); o Macro F1
         acompanha como análise de robustez};

    % ---------------- amarração entre as duas réguas -----------------------
    % conector vertical reto entre as duas caixas de decisão, com o rótulo
    % ancorado à esquerda da linha (nada cruza título nem seta).
    % Fica por último porque referencia nós das duas faixas.
    \draw[seta, dashed] (gate.south) -- (crit.north);
    \node[rot, text width=40mm, anchor=east] at (4.3,-3.3)
        {a mesma razão de $0{,}95$ nas duas réguas:
         $85\% \approx 0{,}95 \times 89{,}56\%$ do melhor
         baseline supervisionado leve};
\end{tikzpicture}
\end{document}
%   (no modo standalone, troque as \ref{...} por texto fixo ou defina-as)
% Requer apenas as bibliotecas já carregadas em packages.tex
% (arrows.meta, positioning). Desenhado para ser legível em preto e branco:
% decisão = borda grossa + cinza escuro; medição = cinza claro;
% consequência = cinza médio; condicional/remissão = tracejado.
% FREEZE respeitado: todos os números são os já fixados no texto
% (85%; 0,95 x 89,56%; 0,95 x acc(D); 34.724; 231.490; 15%).
% Compilado e inspecionado visualmente (pdflatex + render PNG) — ver
% NOTA-esquemas.md, iterações 4-5 do loop.
% ============================================================================
\begin{tikzpicture}[
    x=1cm, y=1cm,
    caixa/.style={draw, rounded corners=2pt, align=center, text width=34mm,
                  minimum height=11mm, inner sep=3pt, font=\footnotesize},
    decisao/.style={caixa, thick, fill=gray!25},
    medida/.style={caixa, fill=gray!8},
    consequencia/.style={caixa, fill=gray!14},
    seta/.style={-{Stealth[length=2.4mm]}, thick},
    rot/.style={font=\scriptsize, align=center, inner sep=2pt}
]
    % ---------------- faixa superior: o GATE do oráculo -------------------
    \node[rot, anchor=west] (tit1) at (-1.9,1.7)
        {\textbf{Gate do oráculo} (fixado de antemão, Seção~\ref{sec:metodo-oraculo-decisao})};

    \node[medida]  (e0)   at (0,0)
        {avaliação dos oráculos\\LLM (S-rand $n=1.000$;\\S-strat $n\approx1.863$)};
    \node[decisao] (gate) at (4.6,0)
        {algum oráculo com\\acurácia $\ge 85\%$\\na S-rand?};
    \node[caixa]   (sim)  at (9.4,1.3)
        {oráculo \textbf{aprovado}:\\papéis pelo critério\\original};
    \node[consequencia] (nao) at (9.4,-1.3)
        {robustez ao ruído\\\textbf{obrigatória}; oráculo ruidoso\\vira cenário central};

    \draw[seta] (e0) -- (gate);
    % âncoras deixam o rótulo inteiro do lado livre da diagonal
    % (a seta cortava o "m"/"o" — achado da inspeção visual)
    \draw[seta] (gate.east) -- node[rot, anchor=south east, pos=0.55] {sim} (sim.west);
    \draw[seta] (gate.east) -- node[rot, anchor=north east, pos=0.55] {não} (nao.west);

    \node[consequencia, text width=44mm] (papeis) at (9.4,-3.5)
        {papéis: LLM Inicial pela razão acurácia/custo; LLM Avançado se
         superior com significância (senão a Fase~3 é eliminada)};
    \draw[seta] (nao) -- (papeis);

    \node[rot, text width=50mm, anchor=north] (ramo) at (9.4,-4.75)
        {o ramo efetivamente seguido, e a divergência que ele exige
         declarar, são reportados com os resultados (Seção~\ref{sec:res-gate})};
    \draw[seta, dashed] (ramo.north) -- (papeis.south);

    % ---------------- faixa inferior: a RÉGUA da hipótese ------------------
    \node[rot, anchor=west] (tit2) at (-1.9,-5.9)
        {\textbf{Critério da hipótese} (Seção~\ref{sec:intro-hipotese})};

    \node[medida] (regua) at (0,-7.4)
        {régua $D$: o mesmo\\classificador com o\\\textit{pool} de referência\\inteiramente rotulado};
    \node[decisao] (crit) at (4.6,-7.4)
        {$\mathrm{acc} \ge 0{,}95 \times \mathrm{acc}(D)$\\com no máximo\\34.724 rótulos?};
    \node[caixa] (veredito) at (9.4,-7.4)
        {veredito da hipótese\\(julgado nos resultados,\\Seção~\ref{sec:res-e3p})};

    \draw[seta] (regua) -- (crit);
    \draw[seta] (crit) -- (veredito);
    \node[rot, below=2.5mm of crit, text width=62mm]
        {teto de 34.724 rótulos $= 15\%$ da população deduplicada de 231.490;
         a acurácia é a métrica do critério (pré-registro); o Macro F1
         acompanha como análise de robustez};

    % ---------------- amarração entre as duas réguas -----------------------
    % conector vertical reto entre as duas caixas de decisão, com o rótulo
    % ancorado à esquerda da linha (nada cruza título nem seta).
    % Fica por último porque referencia nós das duas faixas.
    \draw[seta, dashed] (gate.south) -- (crit.north);
    \node[rot, text width=40mm, anchor=east] at (4.3,-3.3)
        {a mesma razão de $0{,}95$ nas duas réguas:
         $85\% \approx 0{,}95 \times 89{,}56\%$ do melhor
         baseline supervisionado leve};
\end{tikzpicture}
